% ===== §14 =====
\section{The Throwing of a Generative Thing}
\label{p27:sec:throwing}

\noindent
The fifth capacity is the most distinctive of the model and the least reducible to a technique, and it is best approached through what it is not. When one has entered another's world and observed a matter the other has not observed, the impulse is to tell the other what one has seen, to deliver the observation as a conclusion. This delivery, however well meant, forces upon the other a rendering formed in one's own movement, and it places the other in the position of receiving a finished result and does not generate one. Even when the result is right, in the provisional sense the account allows, its delivery forecloses the other's own generation and installs one's own observation in the other's world from without. The fifth capacity is the alternative to this delivery, and it consists in giving the other the means of their own generation and not a conclusion.

To throw a generative thing into the other's unfolding is to offer, in the other's own terms and in keeping with the generative rules of the other's world, something from which the other will generate an observation they had not reached, an observation that is then theirs and not a transplant of one's own. The thing thrown may be a question that opens a direction the other's world contains but has not taken, a juxtaposition that sets two of the other's own renderings into a relation the other had not drawn, an instance that the other's own commitments make significant. Its mark is that it is generative and not conclusive, that it occasions the other's generation of the observation and does not carry it as a payload to be delivered, the other's generation of it, and that what the other then observes is generated in the other's own world and belongs to the other.

This is the capacity that a generative pedagogy has practiced under other names, and the theory of generative justice developed by Eglash~\parencite{eglash2016generative} gives it a wider setting, in which what a generative act returns is returned to those who generate it and is not extracted from them. The bearing of that setting on the present capacity is exact. To throw a generative thing rightly is to act so that the observation generated returns to the other in whom it is generated, remaining the other's own, where the forcing of a conclusion extracts the generation, installing one's own result in place of the other's and taking from the other the generativity that was theirs to exercise. The distinction between the throwing and the forcing is thus a distinction of justice and not of tact, between an act that leaves the other the author of what they come to observe and an act that authors it for them.

Three features distinguish the throwing from its failures, and they can be named though they cannot be reduced to a rule. The first is that the thing thrown must be in the other's terms and must respect the generative rules of the other's world, for a thing thrown in one's own terms is a delivery in disguise, and a thing that violates the other's rules will only collide with the other's world and will not generate within it. This is why the throwing depends upon the prior capacities, the recognition and the entry and the mapping, since one cannot offer in the other's terms what one has not learned to render in them. The second is that the thing thrown must be generative and not conclusive, an opening and not an answer, for a conclusion pressed in the form of a question is a forcing that has borrowed the shape of a throwing, and the other feels the difference as the difference between being drawn out and being led. The third is a matter of moment and measure, of when the thing is offered and how much is offered, for a generative thing thrown before the other's world is ready to receive it will not generate, and one thrown with too heavy a hand will steer where it should open. This third feature cannot be given as a rule at all, and it is the point at which the capacity passes from a skill into the standing disposition the previous part described, since the sense of the moment is formed and not followed.

The sixth capacity belongs with the fifth, as the withholding that the throwing requires once it is done. To throw a generative thing and then to steer the other's generation of it, to guide the other toward the observation one had oneself reached, is to take back with the sixth capacity what the fifth had given, converting the occasion of the other's generation into a managed route to one's own conclusion. Non-intervention is the capacity to let go of the thing thrown, to allow the shared dynamic between the two worlds to develop without one's steering, and to bear the uncertainty of not knowing what the other will generate or whether they will generate at all. It is a capacity of restraint, and its difficulty is the difficulty of trusting a process one has set going but does not control. Its characteristic failure is the intervention that cannot wait, the resumption of steering under the anxiety of an open outcome, which returns the exchange to the delivery the fifth capacity was meant to escape.

Together the fifth and the sixth capacities compose a single gesture in two moments, the giving of a generative thing and the letting go of it, the offering into the other's world and the withholding of one's hand from what the other then makes. The gesture is a giving in the strong sense, since what it gives is the other's own generation of an observation and not an observation itself, and its condition is a restraint that wants nothing back from the giving, neither the credit of the result nor its conformity to what one had oneself observed. This restraint is the form that non-possession takes in the fifth and sixth capacities, and it is what holds the throwing apart from the refined management that would use the appearance of a generative gift to install a conclusion of its own. The disposition to give in this way, wanting nothing back, is not a technique that can be rehearsed to reliability. It is a formed generosity, cultivated as a way of being toward the other and carried into the wager, where its rightness cannot be confirmed at the moment of the giving and must be borne on trust.
